\section*{Bab \ref{ch:g11:binom} --- Distribusi Binomial}

\begin{solution}{exo:g11:binom:1}
Baris $6$ dan $7$:
\[
1,\ 6,\ 15,\ 20,\ 15,\ 6,\ 1
\qquad\text{dan}\qquad
1,\ 7,\ 21,\ 35,\ 35,\ 21,\ 7,\ 1 .
\]
Jadi $\binom62 = 15$, $\binom73 = 35$, $\binom74 = 35$ (kesimetrian
$\binom73 = \binom74$ mencerminkan pertukaran keberhasilan dan kegagalan).
\end{solution}

\begin{solution}{exo:g11:binom:2}
Banyaknya percobaan tetap ($4$ lemparan), dua hasil per lemparan (enam atau
bukan, $p = \frac16$), lemparannya saling bebas: $X \sim \mathcal B(4, \frac16)$.
\[
\P(X=0) = \left(\frac56\right)^4 = \frac{625}{1296} \approx 0.48,
\qquad
\P(X=1) = 4 \times \frac16\left(\frac56\right)^3 = \frac{500}{1296}
\approx 0.39,
\]
\[
\P(X \geq 2) = 1 - \frac{625 + 500}{1296} = \frac{171}{1296}
\approx 0.13 .
\]
\end{solution}

\begin{solution}{exo:g11:binom:3}
\emph{1.} Binomial $\mathcal B(20, \frac12)$: $n$ tetap, $p$ sama,
lemparannya saling bebas.

\emph{2.} Bukan binomial: kartunya dibagikan \emph{tanpa pengembalian},
sehingga peluang memperoleh kartu hati berubah dari kartu ke kartu dan
pengambilannya tidak saling bebas.

\emph{3.} Binomial $\mathcal B(7, 0.3)$ menurut asumsi kebebasan yang
dinyatakan.
\end{solution}

\begin{solution}{exo:g11:binom:4}
$\E(X) = 50 \times 0.2 = 10$;
$\V(X) = 50 \times 0.2 \times 0.8 = 8$;
$\sigma(X) = 2\sqrt2 \approx 2.83$.
\end{solution}

\begin{solution}{exo:g11:binom:5}
$X \sim \mathcal B(6, 0.7)$.
\[
\P(X = 4) = \binom64 (0.7)^4 (0.3)^2 = 15 \times 0.2401 \times 0.09
\approx 0.324 .
\]
\[
\P(X \geq 5) = \binom65 (0.7)^5(0.3) + (0.7)^6
= 6 \times 0.16807 \times 0.3 + 0.117649 \approx 0.420 .
\]
\end{solution}

\begin{solution}{exo:g11:binom:6}
$X \sim \mathcal B\left(8, \frac12\right)$; setiap jalurnya berpeluang
$\frac{1}{256}$, sehingga
\[
\P(X \geq 6) = \frac{\binom86 + \binom87 + \binom88}{256}
= \frac{28 + 8 + 1}{256} = \frac{37}{256} \approx 0.14 .
\]
Menerka lulus kira-kira satu kali dari tujuh.
\end{solution}

\begin{solution}{exo:g11:binom:7}
\omterm{def:g10:proba:operations}{Komplemen} ``paling sedikit satu dari setiap jenis'' adalah ``kelimanya sama'':
semuanya A atau semuanya B, masing-masing berpeluang $\left(\frac12\right)^5 =
\frac1{32}$. Jadi
\[
\P(\text{satu dari setiap jenis}) = 1 - \frac{2}{32} = \frac{15}{16} .
\]
\end{solution}

\begin{solution}{exo:g11:binom:8}
$\P(X = 3) = p^3$ dan $\P(X \geq 1) = 1 - (1-p)^3$. Menyelesaikan
$p^3 = \frac{27}{64} = \left(\frac34\right)^3$ memberi $p = \frac34$ (\omterm{def:g11:func:function}{fungsi}
pangkat tiga tegas naik, \cref{ch:g11:func}, sehingga penyelesaiannya
tunggal).
\end{solution}

\begin{solution}{exo:g11:binom:9}
$\P(\text{paling sedikit satu gambar}) = 1 - \left(\frac12\right)^n$, sehingga
syaratnya adalah $\left(\frac12\right)^n < 0.01$, yaitu $2^n > 100$. Karena
$2^6 = 64$ dan $2^7 = 128$: mulai dari $n = 7$ lemparan.
\end{solution}

\begin{solution}{exo:g11:binom:10}
Jumlah baris $n$ membilang semua jalur pada \omterm{met:g10:proba:tree}{pohon} $n$ percobaan, yang
dikelompokkan menurut banyaknya keberhasilan. Namun \omterm{met:g10:proba:tree}{pohonnya} melipatduakan
jalurnya pada setiap percobaan (setiap jalur terbelah menjadi B dan G),
sehingga seluruhnya ada $2^n$ jalur. Jadi
$\sum_{k} \binom nk = 2^n$. Cara lain, lewat induksi: baris $0$ berjumlah
$1 = 2^0$, dan aturan Pascal membuat setiap entri baris $n$ menyumbang ke
tepat dua entri baris $n+1$, sehingga jumlah barisnya berlipat dua.
\end{solution}

\begin{solution}{exo:g11:binom:11}
\emph{1.} Di bawah pengakuan itu, $X \sim \mathcal B(10, 0.6)$. Menjumlahkan
suku pertamanya:
\begin{align*}
\P(X \leq 3)
&= (0.4)^{10} + 10(0.6)(0.4)^9 + 45(0.6)^2(0.4)^8
+ 120(0.6)^3(0.4)^7\\
&\approx 0.0001 + 0.0016 + 0.0106 + 0.0425 = 0.0548 .
\end{align*}

\emph{2.} Hasil yang paling sedikit seekstrem yang teramati ($3$ pendukung
atau kurang) berpeluang sekitar $5.5\%$ --- tepat \emph{di atas} ambang
$5\%$. Dengan menerapkan aturannya secara ketat, pengamatannya (nyaris)
sesuai dengan pengakuan itu dan kita tidak menolaknya. Contoh ini memperlihatkan
betapa pekanya keputusan di ambang batas terhadap pilihan ambangnya: dengan
kelaziman $6\%$ kesimpulannya akan berbalik.
\end{solution}

\begin{solution}{pb:g11:binom:1}
\textbf{1.} Barisnya: $1$; $1\,1$; $1\,2\,1$; $1\,3\,3\,1$;
$1\,4\,6\,4\,1$; $1\,5\,10\,10\,5\,1$;
$1\,6\,15\,20\,15\,6\,1$. Kesimetriannya: memilih $k$ objek mana
yang diambil adalah tindakan yang sama dengan memilih $n - k$ mana yang
ditinggalkan.

\textbf{2.} $1 + 4 + 6 + 4 + 1 = 16 = 2^4$;
$1 + 5 + 10 + 10 + 5 + 1 = 32 = 2^5$. Buktinya: semua
$\binom nk$ membilang himpunan bagian dari setiap ukuran pada himpunan
beranggota $n$, dan seluruh himpunan bagiannya berjumlah $2^n$ (setiap
anggotanya masuk atau tidak, secara saling bebas).

\textbf{3.} Panitia beranggota $k$ yang dipilih dari $n + 1$ orang, salah
satunya Zoe: panitia tanpa Zoe ada $\binom nk$ (pilihlah seluruh
$k$ dari yang lain); panitia dengan Zoe ada $\binom{n}{k-1}$
(pilihlah $k - 1$ rekannya). Seluruhnya:
$\binom nk + \binom{n}{k-1}$.

\textbf{4.} Segitiganya, baris $7$: $1\,7\,21\,35\,\dots$: jadi $35$.
Rumusnya: $\frac{7 \times 6 \times 5}{3 \times 2 \times 1} =
35$.

\textbf{5.} $1 + 3 + 6 + 10 = 20 = \binom63$. Penerapan berjenjangnya:
$\binom63 = \binom52 + \binom53 = \binom52 + \binom42 +
\binom43 = \binom52 + \binom42 + \binom32 + \binom33$ ---
setiap penerapan aturan Pascal mengupas satu suku dari
tangganya.

\textbf{6.} Banyaknya pantulan $n$ tetap; setiap pantulannya sebuah
\omterm{def:g11:binom:bernoulli}{percobaan Bernoulli} saling bebas dengan $p = \frac12$ yang sama; nomor
kotaknya membilang keberhasilannya (pantulan ke kanan): jadi ketiga
kotak periksa \cref{met:g11:binom:recognize} tercentang:
$\mathcal B\!\left(n, \frac12\right)$.

\textbf{7.} Peluangnya $\frac{1}{16}, \frac{4}{16},
\frac{6}{16}, \frac{4}{16}, \frac{1}{16}$ untuk kotak
$0, \dots, 4$: kotak tengah $2$ yang paling padat.

\textbf{8.} Harapan cacahnya $= 1024 \times
\binom{10}{k}/1024 = \binom{10}{k}$: kotak tengah
$\binom{10}{5} = 252$ bola; kotak $7$: $\binom{10}{7} = 120$;
setiap kotak tepinya: $1$ bola. Pusat yang tinggi lalu menurun
secara setangkup ke tepi yang setipis bisikan: itulah loncengnya.

\textbf{9.} $\E = np = 5$; $V = np(1 - p) = 2.5$;
$\sigma \approx 1.58$. Dalam $2\sigma$: kotak $2$ sampai $8$
membawa
\[
\frac{45 + 120 + 210 + 252 + 210 + 120 + 45}{1024}
= \frac{1002}{1024} \approx 98\,\%
\]
dari bolanya --- jauh lebih baik
daripada $75\,\%$ serbaguna milik Chebyshev
(\cref{pb:g11:stat:1}): bentuk lonceng memusat dengan kuat.

\textbf{10.} $\E = 6$, $V = 10 \times 0.6 \times 0.4 = 2.4$,
$\sigma \approx 1.55$: tumpukannya mempertahankan bentuk loncengnya tetapi
menggeser puncaknya ke kotak $6$ --- papan yang miring adalah koin yang
berat sebelah, yang dibuat terlihat.

\textbf{11.} Nomor kotaknya adalah \emph{jumlah} banyak dorongan
kebetulan yang kecil, saling bebas, dan sama besarnya --- dan jumlah semacam
itu selalu menata dirinya menjadi lonceng: kebanyakan dorongannya saling
menghapus, sedangkan yang ekstrem menuntut kebulatan suara. Tinggi badan,
galat pengukuran dan tak terbilang besaran alam pun sama-sama jumlah banyak
pengaruh kecil yang saling bebas, dan itulah sebabnya bayangan yang sama
muncul di mana-mana; teorema yang menjaminnya adalah teorema limit pusat.

\textbf{12.} Sapu bersih: satu tim memenangkan semua $4$ laganya:
$2 \times \left(\frac12\right)^4 = \frac18$.

\textbf{13.} Tujuh laga menuntut kedudukan $3$--$3$ setelah enam laga:
$\binom63 \left(\frac12\right)^6 = \frac{20}{64} =
\frac{5}{16}$.

\textbf{14.} Berakhir pada laga ke-$5$: pemenangnya merebut laga 5 dan $3$
dari $4$ laga pertamanya:
$2 \times \binom43 \left(\frac12\right)^5 = \frac14$. Berakhir pada laga
ke-$6$: $2 \times \binom53 \left(\frac12\right)^6 = \frac{5}{16}$.
\omterm{def:g11:prob:rv}{Distribusinya} atas $4, 5, 6, 7$:
$\frac18, \frac14, \frac{5}{16}, \frac{5}{16}$ (berjumlah $1$).
Harapan panjangnya: $4 \cdot \frac18 + 5 \cdot \frac14 + 6 \cdot
\frac{5}{16} + 7 \cdot \frac{5}{16} = 5.8125$ laga. Rangkaian enam dan
tujuh laga yang paling mungkin --- ketegangannya sudah terpasang di dalam
formatnya.

\textbf{15.} Menang dalam $4$: $0.6^4 = 0.1296$; dalam $5$:
$\binom43\,0.6^3 \times 0.4 \times 0.6 = 0.2074$; dalam $6$:
$\binom53\,0.6^3 \times 0.4^2 \times 0.6 = 0.2074$; dalam $7$:
$\binom63\,0.6^3 \times 0.4^3 \times 0.6 = 0.1659$. Seluruhnya:
sekitar $0.710$: jadi tim berpeluang $60\,\%$ per laga memenangkan
$71\,\%$ rangkaian --- rangkaiannya memperbesar keunggulannya.

\textbf{16.} Final tunggal: $60\,\%$. Sistem tiga laga:
$p^2 + 2p^2 q = 0.36 + 0.288 = 0.648$. Tangga
$60\,\% \to 65\,\% \to 71\,\%$ berlanjut seiring panjangnya: makin banyak
laga makin merata-ratakan keberuntungan (hukum bilangan besar dalam ukuran
kecil), sehingga final yang panjang menobatkan keterampilan --- dan itulah
yang justru dijual oleh liga.

\textbf{17.} Koin setimbang: $\E = 50$, $\sigma = \sqrt{25} = 5$;
$z = \frac{62 - 50}{5} = 2.4$: melampaui kelaziman $2\sigma$
--- jadi koin itu layak diselidiki.

\textbf{18.} $\P(X = 0) = 0.98^{50} \approx 0.364$;
$\P(X = 1) = 50 \times 0.02 \times 0.98^{49} \approx 0.372$;
$\P(X = 2) = \binom{50}{2} 0.02^2 \times 0.98^{48} \approx
0.186$. Jadi $\P(X \geq 3) \approx 1 - 0.922 = 0.078$: sekitar
$7.8\,\%$ --- di atas ambang $5\,\%$, sehingga belum
cukup untuk menolak pengakuannya; kumpulan buruk yang kedua akan mengubah
ceritanya.

\textbf{19.} $\P(\text{paling sedikit satu menang}) = 1 - 0.999^{1000}
\approx 1 - 0.368 = 0.632$: seribu karcis dengan peluang satu per
seribu memberi bukan kepastian melainkan $63\,\%$. Bilangan
$0.368$ yang berulang itu adalah $\frac1e$ dalam samaran --- tetapan $e$
membuat penampilan resminya pada tahun berikutnya.

\textbf{20.} Pengenalannya: $n$ tetap, kebebasan, $p$ yang
tetap --- baru setelah itu, binomial. Tabelnya: \omterm{prop:g11:binom:pascal}{segitiga
Pascal}, baris $n$. Bentuknya: lonceng, setangkup untuk
$p = \frac12$, tergeser bila selain itu. Kompasnya: pusat $np$,
sebaran $\sqrt{np(1-p)}$ --- yaitu \omterm{pb:g11:stat:1}{skor-z} pada setiap keputusan.
Pewarisnya: kurva lonceng mulus yang didekati tumpukannya, dan
hukum bilangan besar yang menjelaskan mengapa papan besar tidak pernah berdusta.

\end{solution}
